
\documentclass{article}
\usepackage{colortbl}
\usepackage{makecell}
\usepackage{multirow}
\usepackage{supertabular}

\begin{document}

\newcounter{utterance}

\centering \large Interaction Transcript for game `cladder', experiment `full\_v1.5\_default', episode 8830 with Qwen3.5{-}4B.
\vspace{24pt}

{ \footnotesize  \setcounter{utterance}{1}
\setlength{\tabcolsep}{0pt}
\begin{supertabular}{c@{$\;$}|p{.15\linewidth}@{}p{.15\linewidth}p{.15\linewidth}p{.15\linewidth}p{.15\linewidth}p{.15\linewidth}}
    \# & \multicolumn{2}{c}{Player} && \multicolumn{2}{c}{Game Master} \\
    \hline

    \theutterance \stepcounter{utterance}  
    & & & \multicolumn{4}{p{0.6\linewidth}}{
        \cellcolor[rgb]{0.9,0.9,0.9}{
            \makecell[{{p{\linewidth}}}]{
                \texttt{\tiny{[P1$\langle$GM]}}
                \texttt{You are an expert in causal inference. The following question is not a typical commonsense query, but rather a meticulously designed question created by a professor specializing in causal inference, intended to assess the students' mastery of the course content.} \\
\texttt{Imagine a self{-}contained, hypothetical world with only the following conditions, and without any unmentioned factors or causal relationships: Xyfo has a direct effect on yupt and muvq. Zupj has a direct effect on yupt. Yupt has a direct effect on muvq. Xyfo is unobserved. Method 1: We look at how yupt correlates with muvq case by case according to xyfo. Method 2: We look directly at how yupt correlates with muvq in general. To understand how yupt affects muvq, is it more correct to use the Method 1 than Method 2?} \\
\texttt{Start your answer with 'Yes' or 'No', followed by additional reasoning or evidence to support your explanation.} \\
            }
        }
    }
     \\ \\

    \theutterance \stepcounter{utterance}  
    & & \multicolumn{4}{p{0.6\linewidth}}{
        \cellcolor[rgb]{0.95,0.95,0.95}{
            \makecell[{{p{\linewidth}}}]{
                \texttt{\tiny{[P1$\leftrightarrow$P]}}
                \texttt{{[}DRAFT{]}} \\
\texttt{Yes} \\
\\ 
\texttt{**Reasoning:**} \\
\\ 
\texttt{In this hypothetical world, the causal structure is defined as follows:} \\
\texttt{1.  \$X\_\{yfo\} \textbackslash{}rightarrow Y\_\{upt\}\$} \\
\texttt{2.  \$Z\_\{upj\} \textbackslash{}rightarrow Y\_\{upt\}\$} \\
\texttt{3.  \$Y\_\{upt\} \textbackslash{}rightarrow M\_\{uvq\}\$} \\
\\ 
\texttt{This forms a **collider** structure at \$Y\_\{upt\}\$ (specifically, \$X\_\{yfo\} \textbackslash{}rightarrow Y\_\{upt\} \textbackslash{}leftarrow Z\_\{upj\}\$) if we consider the joint distribution of \$X\_\{yfo\}\$ and \$Z\_\{upj\}\$, but more importantly, \$Y\_\{upt\}\$ acts as a **mediator** between the unobserved cause \$X\_\{yfo\}\$ and the outcome \$M\_\{uvq\}\$.} \\
\\ 
\texttt{Here is why Method 1 is superior to Method 2 for understanding the specific effect of \$X\_\{yfo\}\$ on \$M\_\{uvq\}\$:} \\
\\ 
\texttt{1.  **The Problem with Method 2 (General Correlation):**} \\
\texttt{    Method 2 looks at the marginal correlation between \$Y\_\{upt\}\$ and \$M\_\{uvq\}\$ without conditioning on \$X\_\{yfo\}\$. While \$Y\_\{upt\}\$ does causally affect \$M\_\{uvq\}\$, the observed correlation in the general population is confounded by the unobserved variable \$X\_\{yfo\}\$. Because \$X\_\{yfo\}\$ causes \$Y\_\{upt\}\$ and \$Y\_\{upt\}\$ causes \$M\_\{uvq\}\$, there is a valid causal path. However, if \$X\_\{yfo\}\$ is correlated with other factors affecting \$Y\_\{upt\}\$ (like \$Z\_\{upj\}\$) or if we are trying to isolate the *specific* contribution of \$X\_\{yfo\}\$, the marginal correlation mixes the effects of all causes of \$Y\_\{upt\}\$. More critically, if we were trying to estimate the causal effect of \$X\_\{yfo\}\$ on \$M\_\{uvq\}\$ using only \$Y\_\{upt\}\$ and \$M\_\{uvq\}\$ without knowing \$X\_\{yfo\}\$, we cannot separate the influence of \$X\_\{yfo\}\$ from \$Z\_\{upj\}\$ or other unobserved confounders of \$Y\_\{upt\}\$.} \\
\\ 
\texttt{2.  **The Advantage of Method 1 (Conditioning on the Cause):**} \\
\texttt{    Method 1 involves looking at the relationship between \$Y\_\{upt\}\$ and \$M\_\{uvq\}\$ *case by case according to \$X\_\{yfo\}\$*. This is equivalent to conditioning on the value of the unobserved cause \$X\_\{yfo\}\$.} \\
\texttt{    *   By fixing \$X\_\{yfo\}\$, we block the flow of variation from \$X\_\{yfo\}\$ to \$Y\_\{upt\}\$.} \\
\texttt{    *   In this conditioned stratum, the variation in \$Y\_\{upt\}\$ is driven solely} \\
            }
        }
    }
    & \\ \\

    \theutterance \stepcounter{utterance}  
    & & \multicolumn{4}{p{0.6\linewidth}}{
        \cellcolor[rgb]{0.95,0.95,0.95}{
            \makecell[{{p{\linewidth}}}]{
                \texttt{\tiny{[P1$\leftrightarrow$P]}}
                \texttt{{[}SELF{-}FEEDBACK{]}} \\
\texttt{FORMAT: problem} \\
\texttt{RULES: problem} \\
\texttt{STRATEGY: problem} \\
\texttt{REASON: The answer is incomplete, fails to address the specific causal inference question about unobserved confounding, and contains a fundamental logical error regarding collider bias.} \\
            }
        }
    }
    & \\ \\

    \theutterance \stepcounter{utterance}  
    & \multicolumn{4}{p{0.6\linewidth}}{
        \cellcolor[rgb]{0.561,0.737,0.561}{
            \makecell[{{p{\linewidth}}}]{
                \texttt{\tiny{[P1$\rangle$GM]}}
                \texttt{Yes} \\
\\ 
\texttt{**Reasoning:**} \\
\\ 
\texttt{In this hypothetical world, the causal structure is defined as follows:} \\
\texttt{1.  \$X\_\{yfo\} \textbackslash{}rightarrow Y\_\{upt\}\$} \\
\texttt{2.  \$Z\_\{upj\} \textbackslash{}rightarrow Y\_\{upt\}\$} \\
\texttt{3.  \$Y\_\{upt\} \textbackslash{}rightarrow M\_\{uvq\}\$} \\
\\ 
\texttt{Here, \$Y\_\{upt\}\$ acts as a **mediator** between the causes \$X\_\{yfo\}\$ and \$Z\_\{upj\}\$ and the outcome \$M\_\{uvq\}\$. The question asks whether it is more correct to use Method 1 (conditioning on \$X\_\{yfo\}\$) than Method 2 (marginal correlation) to understand how \$X\_\{yfo\}\$ affects \$M\_\{uvq\}\$.} \\
\\ 
\texttt{The answer is **Yes**, and here is why:} \\
\\ 
\texttt{1.  **The Problem with Method 2 (Marginal Correlation):**} \\
\texttt{    Method 2 examines the total correlation between \$Y\_\{upt\}\$ and \$M\_\{uvq\}\$ across the entire population. This correlation aggregates the effects of all causes of \$Y\_\{upt\}\$, including both the unobserved \$X\_\{yfo\}\$ and the observed \$Z\_\{upj\}\$. Because \$X\_\{yfo\}\$ is unobserved, we cannot directly measure its specific contribution to the variation in \$Y\_\{upt\}\$. Consequently, the marginal correlation between \$Y\_\{upt\}\$ and \$M\_\{uvq\}\$ is a composite signal that mixes the causal pathway from \$X\_\{yfo\}\$ with the pathway from \$Z\_\{upj\}\$. Without knowing the specific value of \$X\_\{yfo\}\$, we cannot disentangle the unique effect of \$X\_\{yfo\}\$ from the effect of \$Z\_\{upj\}\$ or other unobserved factors.} \\
\\ 
\texttt{2.  **The Advantage of Method 1 (Conditioning on the Cause):**} \\
\texttt{    Method 1 involves analyzing the relationship between \$Y\_\{upt\}\$ and \$M\_\{uvq\}\$ within strata defined by the values of \$X\_\{yfo\}\$. Although \$X\_\{yfo\}\$ is unobserved, the logic of causal inference dictates that to isolate the effect of a specific cause, one must condition on that cause.} \\
\texttt{    *   By fixing \$X\_\{yfo\}\$ to a specific value, we effectively block the variation in \$Y\_\{upt\}\$ that is attributable to \$X\_\{yfo\}\$.} \\
\texttt{    *   Within each stratum of \$X\_\{yfo\}\$, the remaining variation in \$Y\_\{upt\}\$ is driven solely by \$Z\_\{upj\}\$ (and potentially other unobserved factors common to the stratum, but distinct from the variation caused by \$X\_\{yfo\}\$).} \\
\texttt{    *   Therefore, the correlation observed in Method 1 for a fixed \$X\_\{yfo\}\$ represents the effect of \$Y\_\{upt\}\$ on \$M\_\{uvq} \\
            }
        }
    }
    & & \\ \\

    \theutterance \stepcounter{utterance}  
    & & & \multicolumn{2}{p{0.3\linewidth}}{
        \cellcolor[rgb]{0.9,0.9,0.9}{
            \makecell[{{p{\linewidth}}}]{
                \texttt{\tiny{[GM$|$GM]}}
                \texttt{Yes} \\
            }
        }
    }
    & & \\ \\

    \theutterance \stepcounter{utterance}  
    & & & \multicolumn{2}{p{0.3\linewidth}}{
        \cellcolor[rgb]{0.9,0.9,0.9}{
            \makecell[{{p{\linewidth}}}]{
                \texttt{\tiny{[GM$|$GM]}}
                \texttt{yes} \\
            }
        }
    }
    & & \\ \\

    \theutterance \stepcounter{utterance}  
    & & & \multicolumn{2}{p{0.3\linewidth}}{
        \cellcolor[rgb]{0.9,0.9,0.9}{
            \makecell[{{p{\linewidth}}}]{
                \texttt{\tiny{[GM$|$GM]}}
                \texttt{game\_result = WIN} \\
            }
        }
    }
    & & \\ \\

\end{supertabular}
}

\end{document}
